% TODO make hw stylesheet
\documentclass[10pt]{article}
\usepackage[letterpaper,top=1in]{geometry}
\usepackage[dvipsnames,svgnames]{xcolor}
\usepackage[shortlabels]{enumitem}
\usepackage[framemethod=TikZ]{mdframed}
\usepackage{amsmath,amssymb,amsthm}
\usepackage[colorlinks]{hyperref}
\usepackage{multicol}
\usepackage{tabularx}

\hypersetup{
	colorlinks=true,
	linkcolor=blue,
	filecolor=magenta,
	urlcolor=magenta,
	pdftitle={MAT 528 HW}
}

\usepackage{mathtools}
\usepackage{thmtools}
\usepackage[textsize=footnotesize]{todonotes}
\usepackage{titling}

\reversemarginpar
\mdfdefinestyle{mdbluebox}{roundcorner=10pt,innerbottommargin=9pt,
	linecolor=blue,backgroundcolor=TealBlue!5,}
\declaretheoremstyle[headfont=\sffamily\bfseries\color{MidnightBlue},
mdframed={style=mdbluebox},]{thmbluebox}

\mdfdefinestyle{mdbloobox}{frametitlefont=\bfseries,innerbottommargin=8pt,
	nobreak=true,backgroundcolor=TealBlue!5,linecolor=blue,}
\declaretheoremstyle[headfont=\bfseries\color{MidnightBlue},
mdframed={style=mdbloobox},headpunct={\\[3pt]},postheadspace=0pt,]{thmbloobox}

\mdfdefinestyle{mdredbox}{frametitlefont=\bfseries,innerbottommargin=8pt,
	nobreak=true,backgroundcolor=Salmon!5,linecolor=RawSienna,}
\declaretheoremstyle[headfont=\bfseries\color{RawSienna},
mdframed={style=mdredbox},headpunct={\\[3pt]},postheadspace=0pt,]{thmredbox}

\mdfdefinestyle{mdgreenbox}{linecolor=ForestGreen,backgroundcolor=ForestGreen!5,
	linewidth=2pt,rightline=false,leftline=true,topline=false,bottomline=false,}
\declaretheoremstyle[headfont=\bfseries\sffamily\color{ForestGreen!70!black},
mdframed={style=mdgreenbox},headpunct={ --- },]{thmgreenbox}

\mdfdefinestyle{mdorangebox}{linecolor=RedOrange,backgroundcolor=RedOrange!5,
	linewidth=2pt,rightline=false,leftline=true,topline=false,bottomline=false,}
\declaretheoremstyle[headfont=\bfseries\sffamily\color{RedOrange!70!black},
mdframed={style=mdorangebox},headpunct={ --- },]{thmorangebox}

\mdfdefinestyle{mdblackbox}{linecolor=black,backgroundcolor=RedViolet!5!gray!5,
	linewidth=3pt,rightline=false,leftline=true,topline=false,bottomline=false,}
\declaretheoremstyle[mdframed={style=mdblackbox}]{thmblackbox}

\declaretheoremstyle[spaceabove=6pt,spacebelow=6pt]{basehead}

\declaretheorem[style=basehead,name=Problem]{problem}
\declaretheorem[style=thmredbox,name=Problem,numbered=no]{problem*}
\declaretheorem[style=thmbloobox,name=Setup]{setup} % for config geo
\declaretheorem[style=thmredbox,name=Example,sibling=problem]{example}
\declaretheorem[style=thmredbox,name=$\bigstar$ Problem,sibling=problem]{niceprob}
\declaretheorem[style=thmbluebox,name=Theorem,sibling=problem]{theorem}
\declaretheorem[style=thmbluebox,name=Corollary,sibling=theorem]{corollary}
\declaretheorem[style=thmbluebox,name=Proposition,sibling=theorem]{prop}
\declaretheorem[style=thmbluebox,name=Lemma,numberwithin=problem]{lemma}
\declaretheorem[style=thmbluebox,name=Theorem,numbered=no]{theorem*}
\declaretheorem[style=thmbluebox,name=Lemma,numbered=no]{lemma*}
\declaretheorem[style=thmgreenbox,name=Claim,numberwithin=problem]{claim}
\declaretheorem[style=thmgreenbox,name=Claim,numbered=no]{claim*}
\declaretheorem[style=thmblackbox,name=Remark,numberwithin=problem]{remark}
\declaretheorem[style=thmblackbox,name=Remark,numbered=no]{remark*}
\declaretheorem[style=thmgreenbox,name=Definition,sibling=theorem]{definition}
\declaretheorem[style=thmgreenbox,name=Definition,numbered=no]{definition*}

\providecommand{\ol}{\overline}
\providecommand{\ceil}[1]{\left\lceil #1 \right\rceil}
\providecommand{\floor}[1]{\left\lfloor #1 \right\rfloor}
\newcommand{\dd}{\mathrm{d}}
\providecommand{\ul}{\underline}
\providecommand{\eps}{\varepsilon}
\providecommand{\half}{\frac{1}{2}}
\providecommand{\CC}{\mathbb C}
\providecommand{\FF}{\mathbb F}
\providecommand{\NN}{\mathbb N}
\providecommand{\QQ}{\mathbb Q}
\providecommand{\RR}{\mathbb R}
\providecommand{\ZZ}{\mathbb Z}
\providecommand{\dg}{^\circ}
\providecommand{\ii}{\item}
\providecommand{\setdiff}{\smallsetminus}
\providecommand{\alert}[1]{{\sffamily\textbf{\textcolor{blue}{#1}}}}
\providecommand{\inv}{^{-1}}
\providecommand{\mb}[1]{\mathbf{#1}}

\newenvironment{soln}{\begin{proof}[Solution]}{\end{proof}}

\providecommand{\pow}{\operatorname{Pow}}
\providecommand{\vocab}[1]{{\textbf{\textcolor{ForestGreen}{#1}}}}
\providecommand{\lcm}{\operatorname{lcm}}
\providecommand{\abs}[1]{\left|#1\right|}

\usepackage{mathrsfs}

% place title at top right
\setlength{\droptitle}{-8ex}
\pretitle{\begin{flushright}\Large\bfseries}
\posttitle{\par\end{flushright}}
\preauthor{\begin{flushright}}
\postauthor{\end{flushright}}
\predate{\begin{flushright}}
\postdate{\end{flushright}}

\title{MAT 528 HW07}
\author{\large Aditya Pahuja}
\date{\large Due 10/16}
\begin{document}
\maketitle

\begin{problem}
    [Munkres 28.3]
    Let $X$ be limit point compact.
    \begin{enumerate}[(a)]
        \ii If $f\colon X\to Y$ is continuous,
	does it follow that $f(X)$ is limit point compact?
	\ii If $A$ is a closed subset of $X$, does it follow that
	$A$ is limit point compact?
	\ii If $X$ is a subspace of the Hausdorff space $Z$,
	does it follow that $X$ is closed in $Z$?
    \end{enumerate}
\end{problem}

\begin{soln}
    (a) \framebox{No}.
    Let $Y$ be a two-point set with the indiscrete topology
    and let $X = \NN\times Y$, as in Section 28 Example 1.
    Then, as Munkres shows in the example, $X$ is limit point compact,
    but its image under the projection $\pi_1\colon X\to\NN$
    is $\NN$, which is not limit point compact because it is infinite but
    every point is an isolated point and so no subset has a limit point.
    \\

    (b) \framebox{Yes}.
    If $E\subseteq A$ is infinite, then it is an infinite subset of $X$,
    so $E$ has a limit point $x$. This point $x$ is a limit point of $A$ as well,
    which means $x\in A$ because closed sets contain their limit points,
    which means that $E$ has a limit point in $A$;
    this proves that $A$ is limit point compact.
    \\

    (c) \framebox{No}. Let $\omega_1$ be the first uncountable ordinal
    ($S_{\Omega}$ in Munkres's language, but I will use the standard notation).
    Then $\omega_1+1 = [0,\omega_1]$ is Hausdorff in the order topology,
    and the subspace $[0,\omega_1) = \omega_1$ is limit point compact
    because it is sequentially compact:
    given countable ordinals $\alpha_1$, $\alpha_2$, \dots,
    their supremum (union) $\alpha = \bigcup_{n\in\NN}\alpha_n$ is a countable ordinal
    because unions of countably many countable sets are countable,
    so $\alpha\in\omega_1$, and $\omega$ is a limit point of the sequence
    because if there were an open interval around $\alpha$ not containing
    any $\alpha_n$ then we would be able to take the smaller endpoint of the interval
    to get a smaller upper bound for the sequence.
    To show that $\omega_1$ is not closed in $\omega_1+1$,
    note first that $\omega_1+1$ is compact by Munkres's
    Theorem 27.1, where $[0,\omega_1]$ has the least upper bound property
    because the least upper bound of a set of ordinals is their union.
    On the other hand $[0,\omega_1)$ is not compact because
    the sets of the form $[0,\alpha)$ where $\alpha\in\omega_1$
    form an open cover with no finite subcover.
    Closed subspaces of compact spaces are compact,
    so if $[0,\omega_1)$ were closed then it would be compact;
    therefore it is not closed in $[0,\omega_1]$.
\end{soln}

\begin{problem}
    [Munkres 28.6]
    Let $(X,d)$ be a metric space. If $f\colon X\to X$ satisfies the condition
    $d(f(x), f(y)) = d(x,y)$ for all $x,y\in X$, then
    $f$ is called an isometry of $X$. Show that if $f$ is an isometry
    and $X$ is compact, then $f$ is bijective and hence a homeomorphism.
\end{problem}

\begin{soln}
    $f$ is injective because
    \begin{equation*}
        f(x) = f(y)\implies 0 = d(f(x),f(y)) = d(x,y) \implies x = y.
    \end{equation*}
    Now suppose that $f$ is not surjective, and let $a\in X\setminus f(X)$.
    Since for any $\eps>0$, $d(x,y) < \eps$ implies $d(f(x),f(y)) < \eps$,
    $f$ is continuous, so $f(X)$ is compact, hence closed in $X$
    (since the metric space $X$ is Hausdorff). This means
    we can find $\eps>0$ such that $B_\eps(a)$ is disjoint from $f(X)$.
    Then, define the sequence $x_0 = a$ and $x_n = f(x_{n-1})$ for $n\in\NN$.
    Define $f^1 = f$ and $f^{n+1} = f^n\circ f$ for all $n\in\NN$.
    For any positive integers $m > n$, we have
    \begin{equation*}
	d(x_m, x_n) = d(f^{n}(f^{m-n}(a)), f^n(a))
	= d(f^{m-n}(a), a) \ge \eps,
    \end{equation*}
    where the second equality comes from $f$, hence the iterated function $f^n$,
    being an isometry, while the inequality comes from $f^{m-n}(a)$
    being in $f(X)$ (so it cannot lie in $B_\eps(a)$ by assumption).
    However, since $X$ is compact, the sequence $x_1$, $x_2$, \dots
    has a convergent subsequence, which is a Cauchy sequence,
    so there exist two points in the subsequence
    that are less than $\eps$ away from each other,
    contradicting the inequality that we just proved.
    Therefore, $f$ must be surjective, so together with injectivity
    we have that $f$ is a bijection on $X$.

    We also already established that $f$ (hence any isometry) is continuous,
    and $f\inv$ must therefore be continuous
    because it is also an isometry, since if $x,y\in X$ then
    \begin{equation*}
        d(x,y) = d(f(f\inv(x)), f(f\inv(y))) = d(f\inv(x), f\inv(y))
    \end{equation*}
    where the second equality is due to $f$ being an isometry.
    This proves that $f$ is a continuous bijection with continuous inverse,
    so it is a homeomorphism of $X$.
\end{soln}

\begin{problem}
    [Munkres 28.7]
    Let $(X,d)$ be a metric space. If $f\colon X\to Y$ satisfies the condition
    $d(f(x),f(y)) < d(x,y)$ for all $x,y\in X$ with $x\ne y$,
    then $f$ is called a \vocab{shrinking map}.
    If there is a number $\alpha < 1$ such that
    \begin{equation*}
        d(f(x),f(y)) \le \alpha d(x,y)
    \end{equation*}
    for all $x,y\in X$, then $f$ is a \vocab{contraction}.
    \begin{enumerate}[(a)]
        \ii If $f$ is a contraction and $X$ is compact,
	show that $f$ has a unique fixed point.
	\ii Show more generally that if $f$ is a shrinking map and $X$ is compact,
	then $f$ has a unique fixed point.
	\ii Let $X = [0,1]$. Show that $f(x) = x - x^2/2$
	maps $X$ into $X$ and is a shrinking map that is not a contraction.
	\ii Show that the map $f\colon\RR\to\RR$ given by
	$f(x) = (x + \sqrt{x^2 + 1})/2$ is a shrinking map
	that is not a contraction and has no fixed point.
    \end{enumerate}
\end{problem}

\begin{soln}
    (a) Let $A_n = f^n(X) = \{f^n(x) : x\in X\}$, with $A_0 = X$.
    If $y\in A_n$ then $y = f^n(x) = f^{n-1}(f(x))$ for some $x\in X$
    so $y\in A_{n-1}$ as well, meaning $A_{n-1}\subseteq A_n$
    for all $n\in\NN$. Since $X$ is compact, it is bounded in the metric $d$,
    so there exists some $C>0$ such that $d(x,y) < C$ for all $x,y\in X$.
    Then, we can see inductively that the diameter of $A_n$ is less than $\alpha^nC$:
    the base case $A_0$ is that the diameter of $X$ is less than $C$
    and, if $A_n$ has diameter less than $\alpha^nC$, then
    for any $x,y\in A_{n+1}$, if $x = f(x')$ and $y = f(y')$ for $x',y'\in A_n$,
    \begin{equation*}
	d(x,y) = d(f(x'), f(y')) \le \alpha d(x',y') < \alpha\cdot\alpha^nC
	= \alpha^{n+1}C.
    \end{equation*}

    Since $X$ is compact and each $A_n$ is compact, hence closed,
    because it is the image of compact set $X$,
    this implies that $A = \bigcap_{n=1}^\infty A_n$ is nonempty.
    The diameter of $A$ is bounded above by $\alpha^nC$ for each positive integer $n$,
    so $A$ has diameter zero, meaning it consists of a single point $p$.
    Any fixed point of $f$ must be contained in $A$ because
    if $x\in A_n$ then $f(x) = x\in A_{n+1}$, so this shows that
    if $p$ is a fixed point is is unique.
    To see that $p$ is indeed a fixed point,
    note that $f(A)$ is a subset of $\bigcap_{n=1}^\infty f(A_n) =
    \bigcap_{n=2}^\infty A_n = A$, and $f(A)$ must be nonempty
    because $A$ is nonempty, so $f(A)$ is equal to $A$,
    the only nonempty subset of $A$, meaning $f$ fixes $p$.
    \\

    (b) Define $A_n$ and $A$ as above.
    As in the proof above, $f(A)$ is a subset of $A$.
    We show that $A$ is a subset of $f(A)$. To do this let
    $a\in A$, and for each $n$ pick $x_n\in X$ such that
    $f^{n+1}(x_n) = a$; such $x_n$ exists because $a\in A_{n+1}$.
    Then, the sequence $y_n = f^n(x_n)$, \dots has a convergent subsequence
    $(y_{n_k})_{k=1}^\infty$, say with limit $y$,
    because $X$ is compact, hence sequentially compact.
    If $f$ is a shrinking map it must be continuous
    because for any $\eps>0$, $d(x,y) < \eps$ implies
    $d(f(x), f(y)) < d(x,y) < \eps$,
    and continuity in a metric space is the same thing as sequential continuity, so
    \begin{equation*}
	a = \lim_{n\to\infty} f^{n+1}(x_n) =
	\lim_{k\to\infty} f(y_{n_k}) = f(\lim_{k\to\infty}y_{n_k}) = f(y)
    \end{equation*}
    (since any subsequence of a convergent sequence
    has the same limit as the parent sequence); thus $a\in f(A)$.
    This implies that $A$ and $f(A)$ are equal, as they contain each other.
    Noting that $A$ is compact (because it is the intersection of compact sets),
    $A\times A$ is compact as well, so the continuous function
    $d\colon A\times A\to[0,\infty)$ achieves its maximum
    for some pair of points $(x,y)\in A\times A$; this maximum value is
    by definition the diameter of $A$. Similarly, there exist points
    $x',y'\in f(A)$ such that $d(x',y')$ is the diameter of $f(A)$
    (because $f(A)$ is the continuous image of a compact set, hence compact).
    By proving $A = f(A)$ we have shown that $A$ and $f(A)$ have the same diameter;
    if this diameter were nonzero then, picking $x_0,y_0\in A$ such that
    $f(x_0) = x'$ and $f(y_0) = y'$, we have
    \begin{equation*}
	\operatorname{diam}A \ge d(x_0,y_0) \ge d(f(x_0),f(y_0))
	\ge d(x',y') = \operatorname{diam}f(A),
    \end{equation*}
    so equality must hold in the inequality above;
    in particular $d(x_0,y_0) = d(f(x_0),f(y_0))$ can hold
    if and only if $x_0=y_0$, so that $A = f(A)$ consists of exactly one point;
    the element of $A$ is a fixed point, and its uniqueness follows
    via the same logic as in part (a),
    where any fixed point of $f$ would be contained in all the $A_n$,
    hence also in $A$.
    \\

    (c) $f$ is a shrinking map because for any distinct $x,y\in[0,1]$,
    \begin{equation*}
	\abs{f(x)-f(y)} = \abs{x - y - x^2/2 + y^2/2}
	= \abs{(x-y)(1 - x/2 - y/2)}
	\le \abs{x-y}\cdot\abs{1 - \frac{x+y}2}
    \end{equation*}

    and $0\le \frac{x+y}{2} \le 1$ so $1 - \frac{x+y}2$
    is between $0$ and $1$ as well, hence so is its absolute value; thus
    \begin{equation*}
	\abs{f(x) - f(y)} < \abs{x-y}.
    \end{equation*}
    On the other hand it is not a contraction because
    \begin{equation*}
	1 = f'(0) = \lim_{h\to 0} \frac{f(h) - f(0)}{h},
    \end{equation*}
    so for each $c<1$ there is an $h>0$ such that $f(h) - f(0) > ch$ ---
    to be more precise, if $f$ were a contraction then
    $\sup_{h\in(0,1]} \frac{f(h)-f(0)}{h-0}$ would be less than $1$,
    but the limit above equalling $1$ implies that this supremum
    is at least $1$, so $f$ cnanot be a contraction.
    \\

    (d) $f$ has no fixed point because
    \begin{equation*}
	f(x) = \frac{x + \sqrt{x^2+1}}2 > \frac{x + \sqrt{x^2}}2 \ge \frac{2x}2 = x
    \end{equation*}
    for all $x$. To show that it is a shrinking map, we start by writing
    \begin{equation*}
	\abs{f(x) - f(y)} \le \abs{\frac{x-y}2} +
	\abs{\frac{\sqrt{x^2+1} - \sqrt{y^2+1}}2}.
    \end{equation*}
    Since the derivative of $\frac{\sqrt{t^2+1}}2$ is $\frac{t}{2\sqrt{t^2+1}}$,
    there exists $\xi$ between $x$ and $y$ such that
    the latter addend is equal to
    \begin{equation*}
	\abs{x-y}\cdot \frac{\xi}{2\sqrt{\xi^2+1}}
	< \frac{\abs{x-y}}{2}
    \end{equation*}
    by the mean value theorem; this proves that $\abs{f(x) - f(y)} < \abs{x-y}$.

    To show that $f$ is not a contraction we note that
    \begin{equation*}
	\lim_{x\to\infty}\frac{\abs{f(x)-f(0)}}{\abs{x-0}} =
	\lim_{x\to\infty} \abs{\frac{x + \sqrt{x^2+1} - 1}{2\cdot x}}
	= \lim_{x\to\infty} \abs{\half + \half\sqrt{1 + \frac1{x^2}} - \frac{1}{2x}}
	= 1
    \end{equation*}
    meaning that for any $c<1$, there exists $N_c\in\RR$ such that
    $x > N_c$ implies $\frac{\abs{f(x)-f(0)}}{\abs{x-0}} > c$.
    If $f$ were a contraction then there would exist $\alpha<1$ such that
    $\abs{f(x) - f(0)} \le \alpha\abs{x-0}$, but taking $x > N_\alpha$
    shows that the reverse inequality holds, so $f$ is not a contraction.
\end{soln}

\pagebreak 
\begin{problem}
    [Munkres 29.1]
    Show that the rationals $\QQ$ are not locally compact.
\end{problem}

\begin{soln}
    Let $x\in\QQ$ be an arbitrary point
    and let $C\subseteq\QQ$ be a set containing an open set $U$ around $x$.
    If $C$ were compact, then it would be sequentially compact
    by Munkres's Theorem 28.2. However, we can construct a sequence
    of points in $U$ that does not converge in $C$:
    let $(a,b)\cap\QQ$ be a basis element of the subspace topology on $\QQ$
    that is contained in $U$ and let $r\in (a,b)$ be irrational.
    Then, since $\QQ$ is dense in $\RR$, we can find for each $n\in\NN$
    a rational number $r_n\in (a,b)\cap\QQ$ such that $\abs{r - r_n} < \frac 1n$,
    so that $\lim_{n\to\infty}r_n = r$, which implies any subsequence
    also converges to $r$, so in particular no subsequence of $r_n$
    converges to a point in $\QQ$ (since limits of sequences are unique
    in Hausdorff spaces). Thus there does not exist
    a compact subset of $\QQ$ that contains an open neighborhood of $x$,
    so $\QQ$ is not locally compact.
\end{soln}

\begin{problem}
    Show that the one-point compactification of $\RR$
    is homeomorphic with the circle $S^1$.
\end{problem}

\begin{soln}
    Denote by $\ol\RR$ the one-point compactification of $\RR$
    and let $\omega$ be the point of $\ol\RR\setminus\RR$.
    Define a \vocab{wrapped interval} $)a,b($ as
    $(-\infty,a)\cup(b,\infty)\cup \{\omega\}$, where $a,b$ are real numbers
    such that $a<b$.

    I claim that
    \begin{equation*}
	\mathcal B = \{(a,b) : a,b\in\RR, a<b \} \cup
	\{{)a,b(} : a,b\in\RR, a<b\}
    \end{equation*}
    is a basis for the topology on $\ol\RR$.
    The topology on $\ol\RR$ consists of two kinds of sets:
    \begin{enumerate}[(1)]
        \ii the subsets of $\RR$ that are open in $\RR$; and
	\ii complements in $\ol\RR$ of compact subsets of $\RR$
	(so sets of this kind contain $\omega$).
    \end{enumerate}
    By definition, a wrapped interval is $\ol\RR\setminus[a,b]$
    where $a,b$ are reals with $a<b$, so elements of $\mathcal B$
    are clearly open in $\RR$.
    Conversely, we show that open subsests of $\RR$
    are in the topology generated by $\mathcal B$.
    Sets of the first kind are in the topology generated by $\mathcal B$
    because the open intervlas are in $\mathcal B$.
    To show that sets of the second kind are in $\mathcal B$,
    note that if $C$ is a compact subset of $\RR$, then $C$ is bounded
    by the Heine-Borel theorem, so $C$ can be contained in
    $[a,b]$ for some real numbers $a,b$ with $a<b$;
    picking $\eps>0$ arbitrarily, we can say that
    \begin{equation*}
        \ol\RR\setminus C =
	(\ol\RR\setminus [a,b])\cup
	\underbrace{((a-\eps,b+\eps)\setminus C)}_{U\text{ open in $\RR$}} =
	{)a,b(}\cup U
    \end{equation*}
    so sets of the second kind are unions of open sets in $\RR$
    with a wrapped interval, hence open in the topology generated by $\mathcal B$.

    Let $S^1$ be the unit circle in $\RR^2$ and for $\theta\in\RR$, identify
    the point $(\cos\theta,\sin\theta)\in S^1$ with $\theta$.
    Define $f\colon S^1\to\ol\RR$ by
    \begin{equation*}
        f(0,1) = \omega,\qquad
	f(\cos\theta,\sin\theta) = \frac{\cos\theta}{1-\sin\theta}
	= \tan\left(\frac\theta2 + \frac\pi4\right)
	\text{ if $\sin\theta\ne1$},
    \end{equation*}
    where the expression for $f$ comes from
    \begin{equation*}
        \frac{\cos\theta}{1 - \sin\theta}
	= \frac{\sin(\pi/2-\theta)}{1-\cos(\pi/2-\theta)}
	= \frac{2\sin(\pi/4-\theta/2)\cos(\pi/4-\theta/2)}
	{2\sin^2(\pi/4-\theta)} = \cot\left(\frac\pi4-\frac\theta2\right)
	= \tan\left(\frac\theta2 + \frac\pi4\right).
    \end{equation*}
    Since the value of $f$ increases as $\theta$ increases over $(-3\pi/2,\pi/2)$
    and the restricted map $\widehat f\colon S^1\setminus\{0,1\}\to\RR$ is continuous
    (because it is a quotient of continuous functions),
    $f$ maps open arcs in $S^1$ not containing $\omega$ to open intervals in $\RR$.
    Furthermore, open arcs in $S^1$ containing $(0,1)$
    are really unions of two open arcs, each with one endpoint at $(0,1)$,
    and $\pi/2$. If $3\pi/2 < \alpha < \beta < \pi/2$ are real numbers
    then the counterclockwise arc from
    $(\cos\beta,\sin\beta)$ to $(\cos\alpha,\sin\alpha)$
    contains $(0,1)$, so can be decomposed into
    the counterclockwise arc from $(\cos\beta,\sin\beta)$ to $(0,1)$,
    the counterclockwise arc from $(0,1)$ to $(\cos\alpha,\sin\alpha)$,
    and the point $(0,1)$. Since these two arcs do not contain $(0,1)$
    they get mapped by $f$ to open intervals, where the fact that
    $f$ increases with the angles of the points means that
    the former arc maps to $(f(\cos\beta,\sin\beta),\infty)$
    and the latter arc to $(-\infty,f(\cos\alpha,\sin\alpha))$.
    This means $f$ sends arcs containing $(0,1)$ to wrapped intervals,
    as the two intervals above togetheer with $f(0,1) = \omega$ constitute
    a wrapped interval. This proves that $f$ is an open map,
    as basis elements of the topology on $S^1$ map to
    basis elements of the topology on $\ol\RR$
    so that open sets get sent to open sets.

    Since $\tan(\theta/2 + \pi/4)$ is bijective on $(-3\pi/2,\pi/2)$,
    $\widehat f$ is bijective, so $f$ is also bijective
    (the one extra point $(0,1)$ in the domain of $f$ maps to
    the one extra point $\omega$ in the codomain).

    Let $g\colon(-3\pi/2,\pi/2)\to\RR$ be the map
    $g(\theta) = \tan\left(\frac\theta2 + \frac\pi4\right)$;
    this is bijective. The inverse $g\inv$ is increasing,
    and $f\inv(r) = (\cos(g\inv(r)),\sin(g\inv(r)))$
    when $r\ne\omega$, and of course $f\inv(\omega) = (0,1)$.
    The increasing nature of $g$ means that open intervlas in $\RR$
    map by $f\inv$ to open arcs in $S^1$ not containing $(0,1)$,
    with unbounded open intervals
    mapping to open intervals having $(0,1)$ as an endpoint.
    By decomposing wrapped intervals into unions of intervals with $\{\omega\}$,
    we can see that wrapped intervals map to the open arcs containing $S^1$.
    Thus $f\inv\colon\ol\RR\to S^1$ is also an open map.
    Together with $f$ being an open map, this proves that $f$ is a homeomorphism.
\end{soln}










\end{document}

